\section{随机化给的是概率保证}
\label{sec:08-Numerical-Analysis/05-Numerical-Linear-Algebra/10}

季度审计。审计员拿着两个月前那份报告，要求用同一份数据、同一份代码重跑，核对里面那张主成分表。你照做了。前三个方向几乎一样，第八个方向整体变了号，第十一与第十二互换了位置，解释方差的比例在小数点后第三位对不上。代码没改，数据没改，依赖库的版本也锁着。审计员问：那份报告到底还算不算数？

这不是实现有 bug。上一篇的算法里 \(\Omega\) 是当场抽的，两次运行抽到的不是同一个 \(\Omega\)。那条误差界两次都成立——它保证的从来不是``你会拿到哪个答案''，而是``你拿到的答案有多好''。审计员想核对的是前者，算法承诺的是后者。

这个错位是使用随机化方法时最常见的判断失误，而它与算法本身的正确性无关。

\subsection{两种承诺，形状就不一样}

\hyperref[sec:08-Numerical-Analysis/05-Numerical-Linear-Algebra/01]{``线性方程为何分解而不求逆''}给的是这样一句话：带主元的 LU 分解是后向稳定的，只要增长因子不爆炸，算出的 \(\hat x\) 是某个 \((A+E)\hat x=b\) 的精确解，且 \(\norm{E}\lesssim c(n)\varepsilon\norm{A}\)。这句话对\emph{一切}满足前提的输入成立，没有例外，一次也没有。条件塌掉时它以确定的方式失效——增长因子爆炸，你可以监测它。

上两篇给的是另一句话：对每个固定的输入，随机地抽 \(S\)（或 \(\Omega\)），则\emph{以概率至少 \(1-\delta\)}，结论成立。这里有例外。例外的总测度是 \(\delta\)，而例外落在哪一次运行上，事先无从知道，事后往往也看不出来——失败的那一次不会报错。

\begin{center}
\begin{tikzpicture}[x=1cm,y=1cm,font=\small,>=stealth]
  \node[font=\small,black!80] at (2.3,5.95) {确定性算法};
  \draw[black!50,fill=black!5] (0.2,4.6) rectangle (4.4,5.4);
  \node[font=\footnotesize] at (2.3,5.0) {条件 \(H\) 成立};
  \draw[->,black!55] (2.3,4.55) -- (2.3,4.0);
  \draw[black!50,fill=green!22] (0.2,3.0) rectangle (4.4,3.95);
  \node[font=\footnotesize] at (2.3,3.48) {结论成立};
  \node[align=center,font=\footnotesize,black!70] at (2.3,2.3) {一万次运行\\一万次都成立};
  \node[font=\small,black!80] at (8.9,5.95) {随机化算法};
  \draw[black!50,fill=black!5] (5.8,4.6) rectangle (12.0,5.4);
  \node[font=\footnotesize] at (8.9,5.0) {同一个条件 \(H\) 成立};
  \draw[->,black!55] (8.9,4.55) -- (8.9,4.0);
  \draw[black!50,fill=green!22] (5.8,3.0) rectangle (11.5,3.95);
  \draw[black!50,fill=red!30] (11.5,3.0) rectangle (12.0,3.95);
  \node[font=\footnotesize] at (8.65,3.48) {结论成立，概率 \(1-\delta\)};
  \node[above,font=\scriptsize,red!60!black] at (11.75,3.98) {\(\delta\)};
  \draw[->,black!55] (8.9,2.95) -- (8.9,2.35);
  \node[left,font=\scriptsize,black!70] at (8.75,2.66) {重复 \(t\) 次，留最好的一次};
  \draw[black!50,fill=green!22] (5.8,1.35) rectangle (11.93,2.3);
  \draw[black!50,fill=red!30] (11.93,1.35) rectangle (12.0,2.3);
  \node[font=\footnotesize] at (8.65,1.82) {结论成立，概率 \(1-\delta^{t}\)};
  \node[above,font=\scriptsize,red!60!black] at (11.99,2.33) {\(\delta^{t}\)};
  \node[align=center,font=\footnotesize,black!70] at (8.9,0.65) {失败的那一块压得下去，\\但压不成零};
\end{tikzpicture}

\emph{同样是``条件满足''，两边落下来的东西不一样：左边整条都是结论，右边留了一小块红的，而且事先不知道自己这次落在哪一块。重复几次能把红块压到几乎看不见，却永远压不掉——这不是精度问题，是承诺的类型问题。}
\end{center}

\subsection{失败概率可以压，压的方式要看结果能不能验}

图右下那一步的算术很简单：\(t\) 次独立运行都失败的概率是 \(\delta^t\)。\(\delta=10^{-2}\)、\(t=3\) 就到了 \(10^{-6}\)。这个数量级已经远低于``实现里有 bug''``输入文件被截断''``内存位翻转''这几件事的发生概率，工程上完全够用。

但``留最好的一次''这个动作有个前提，容易被跳过：\emph{你得能便宜地判断哪一次好}。上两篇都满足这个前提。随机化 SVD 里，用几个新的随机向量估 \(\norm{A-QQ^{\trans}A}\) 就能给残差打分，比重跑一次便宜得多；草图最小二乘里，把 \(\tilde x\) 代回原方程算 \(\norm{A\tilde x-b}\) 也是一次矩阵乘向量的事。有廉价的验证，随机化就变成``反复试到满意为止''，失败概率按 \(\delta^t\) 掉。

验证不便宜的时候，这条路走不通。Monte Carlo 估一个没有独立参照的积分，你拿到 \(t\) 个数，没有办法说哪个更接近真值。此时用的是另一招：取 \(t\) 次结果的中位数。只要单次成功的概率超过一半，``过半数的运行同时失败''这件事的概率按 Chernoff 界随 \(t\) 指数下降，结论一样硬，但它换的是另一组前提——要求各次运行独立，且成功的定义是``落在同一个目标区间内''。两种放大手段不能混用。

于是 \(\delta\) 的数值大小基本不是问题。问题是你必须知道自己手上拿的是哪一种保证，以及知道之后该做什么。

\subsection{随机性来自算法，不来自数据}

这一条是本篇最要紧的区分，也是随机化数值算法与统计推断最容易被混为一谈的地方。

统计里的一句典型保证是：若数据独立同分布于某个 \(P\)，则置信区间以概率 \(1-\alpha\) 覆盖真参数。这里的随机性长在\emph{数据}上，保证因而寄生在一个关于数据生成机制的假设上。样本之间有依赖、分布在中途漂移、抽样有选择性——假设一破，那个 \(1-\alpha\) 就不再是 \(1-\alpha\)，而且通常没有任何迹象。\hyperref[sec:11-Mathematical-Statistics/09-Random-Matrix-Theory/03]{``谱偏差怎样影响 PCA 与随机投影''}讲的 PCA 一侧正是这种情形：主方向的估计质量受 \(p/n\) 支配，是数据的性质在决定结论。

上两篇的保证长得不一样：对\emph{任意}的 \(A\) 与 \(b\)，不假设它们服从任何分布、不假设它们之间独立、不假设它们来自任何生成机制，随机地抽 \(S\)，则以概率至少 \(1-\delta\) 结论成立。随机性长在\emph{算法}上，由你自己的随机数发生器提供。

这个差别有一个很实在的后果：没有哪个输入能让随机化算法必然失败。想让它出岔子的人可以针对某一个具体的 \(S\) 构造出坏的 \(A\)，但他不知道你这次会抽到哪个 \(S\)——而 \(S\) 是在看到 \(A\) 之后才抽的。对照\hyperref[sec:08-Numerical-Analysis/05-Numerical-Linear-Algebra/08]{``用随机草图把大矩阵压小''}开头那个均匀抽行的反例就清楚了：那里之所以能构造出必然失败的矩阵，正是因为``均匀抽行''的随机性太弱，输入的结构可以针对它；换成把所有行混起来的随机投影，这条构造路径就断了。

这条推论也有它自己的前提，破起来很轻。它要求 \(S\) 与输入独立。把随机种子写死、公开，并且允许对方按这个种子来准备输入，那么算法在他眼里就是确定性的，最坏情况又回来了。也就是说：\textbf{统计的保证要为输入的性质买单，随机化算法的保证要为随机数的质量买单。}两者的失效方式完全不同，检查清单也完全不同——前者查数据的采集过程，后者查随机源与它和输入之间的独立性。

\subsection{什么时候不该用}

开篇那个审计场景是第一类。需要逐位可复现的结论时，同一份数据两次运行给出不同答案本身就是不合格。补救是固定种子并把种子当作输入的一部分归档。要留神这个动作做了什么：它把算法在操作上变回确定性的，却没有把那条概率保证变成确定性保证。固定种子之后你拿到的是一次特定抽样的结果，好不好由那一次抽样决定了，只是从此不再变。可复现与更准是两件事。

第二类是需要最坏情况保证的场合：安全关键系统、要出具可检验证书的计算、形式验证的一环。这时该走确定性算法；退一步的做法是用随机化快速拿一个候选，再用确定性方法验证或修正——\hyperref[sec:08-Numerical-Analysis/05-Numerical-Linear-Algebra/08]{``用随机草图把大矩阵压小''}里那个先草图后预条件的写法就是这个模式的样板：随机化负责快，确定性负责准。

第三类是矩阵本身不大。\(m,n\) 只有几千时，稠密 LAPACK 例程几秒钟给出确定性答案，而随机化要抽随机数、要做 QR、要估残差、要记录种子，这些固定开销会把省下来的运算量全部吃掉，还额外欠下一份需要解释的报告。规模不到，随机化不划算。

第四类是精度要求高。\(1/\epsilon^2\) 这个依赖关系摆在那里，草图不是拿来求高精度解的工具。

\subsection{三篇合起来说了什么}

三篇用的是同一个动作。\hyperref[sec:08-Numerical-Analysis/05-Numerical-Linear-Algebra/08]{``用随机草图把大矩阵压小''}把随机矩阵放在左边，压掉行数，保住列空间上的长度；\hyperref[sec:08-Numerical-Analysis/05-Numerical-Linear-Algebra/09]{``随机化 SVD：用随机向量探到主子空间''}把它放在右边，压掉列数，保住主子空间；本篇把两者共有的那半句``以高概率''摊在桌面上，看清它究竟承诺了什么。

放回整章看，前七篇每一篇都在拿一个假设换一次运算量的下降：LU 要能选到不太小的主元，Cholesky 要正定，Krylov 要谱分布得当，SVD 要接受秩在浮点下不是整数，Kaczmarz 要系统一致。随机化这三篇换的东西不同——它没有对矩阵提出新要求，它降低的是承诺本身的强度：从``一定''降到``以概率 \(1-\delta\)''。运算量降一个量级，换来的是每一次运行都带着一小块说不清落在哪里的失败区域。

这也是本书那条判断在数值线性代数里的又一次现身：每一层保证的强度都对应一组条件，没有哪种方法凭空变快。随机化并不比 QR 聪明，它只是接受了一个更弱的承诺，并在这个更弱的承诺下把运算量压了下去。用它之前该问的不是``它准不准''，而是``我这件事能不能接受一个概率型的答案''——能接受，它就是当前规模下最合适的工具；不能接受，再快也不该用。

再往下走，被追问的对象会从``解一个方程''变成``跟住一条随时间演化的轨迹''，\hyperref[ch:066]{``常微分方程数值方法：沿局部斜率走完整条轨迹''}接手这件事。保证的形状在那里又会变一次：不再是``这一步算得对不对''，而是``走了一万步之后误差有没有被放大''。
